\section{Background and Related Work}
\label{sec:related}

The introduction argued that characterizing a generator requires path-level
sampling, not only single draws from the terminal marginal. We now summarize the
lines of work that make this program precise: classical TPS and its
extensions to non-reversible dynamics, score-based diffusion and EDM samplers,
structure prediction models and their known failure modes, adaptive biasing
methods including OPES, and symmetry-aware generative modeling on molecular
geometry.

\paragraph{Transition path sampling.}
Classical TPS targets reactive trajectories of (near-)equilibrium molecular
dynamics by Monte Carlo moves in path space, using short-time propagators and,
when available, microscopic reversibility to relate forward and backward
kernels \citep{dellago1998,bolhuis2002}. For dynamics that are not the
time-reversal of a known reference propagator---including learned generators and
fixed-schedule denoisers---\citet{buijsman2020} formulated TPS with explicit
backward proposals on path segments. \citet{falkner2023} conditioned deep
Boltzmann generators for rare events using path weights analogous to those we
use for a frozen denoiser. \citet{falkner2024} revisited shooting moves and
acceptance ratios in this discrete-time, non-reversible setting. Our
development follows these references and instantiates them for a public
protein diffusion pipeline.

\paragraph{Score-based diffusion and EDM.}
Score-based models define a family of noisy laws and learn a score or denoiser
across noise levels \citep{song2021,ho2020}. The EDM parameterization and
sampler design of \citet{karras2022} match engineering practice in many
structure predictors. For TPS, the relevant object is the \emph{discrete}
transition law implemented in code (including optional churn between schedule
levels), not a formal continuous-time limit. Section~\ref{sec:diffusion-boltz}
makes that distinction explicit for Boltz-2.

\paragraph{Generative protein structure models and evaluation.}
Large multimodal predictors have achieved broad chemical coverage
\citep{alphafold3,boltz2,rfaa}. Independent studies have highlighted
miscalibration and memorization risks \citep{gopalan2025,skrinjar2025,hitawala2025}.
Path sampling complements aggregate benchmarks by exposing the generator's
stochastic landscape for a single input conditioning.

\paragraph{Rare events, paths, and other generative modalities.}
Rare-event analysis of autoregressive language models using path-space methods
appears in \citet{dorman2026}. For diffusion on molecular energy landscapes,
\citet{raja2025omtps} connect high-probability paths to an Onsager--Machlup
action under SDE dynamics. Our setting differs: we take the generator's
\emph{discrete} reverse kernel as given and target its induced path measure
without gradient-based path optimization.

\paragraph{Enhanced sampling and OPES.}
Metadynamics and related methods flatten or reshape free-energy barriers along
collective variables \citep{barducci2008}. OPES estimates a marginal along
chosen CVs and applies a well-tempered bias toward a target visit distribution
\citep{invernizzi2020,invernizzi2022}. We use OPES only as an outer bias on the
terminal-frame CVs inside a TPS Metropolis correction
(Section~\ref{sec:opes-theory}), with the caveats of adaptive bias documented
there.

\paragraph{Symmetry, quotients, and equivariance.}
Internal molecular shape is invariant under global rigid motions; quotient-space
diffusion \citep{xu2026quotient} and SE(3) equivariant architectures
\citep{yim2023se,satorras2021,batzner2022,hoogeboom2022} formalize how learning
and sampling should interact with symmetry. Boltz-2 samples in Euclidean
coordinates with random augmentation; we explain how storing augmentation in
the extended TPS state aligns Metropolis weights with the implementation.

Having fixed the literature context, we next give a self-contained derivation of
TPS for abstract discrete-time generators before specializing to diffusion.
